\section{財政検証}

\begin{fillblank}[自作 2026/11/3]
  「令和6(2024)年財政検証結果レポート」第１章の「第１節　国民年金・厚生年金の財政検証の位置づけ　4.財政検証の位置づけ」
  に関する記述について、以下の空欄を埋めなさい。
\\


  我が国の公的年金制度は、高齢者に対する年金の支給に要する費用をそのときの現役世代の負担
  によって賄う「（　　A　　）」を基本としつつ、一定の積立金を保有しそれを活用することにより、
  将来の受給世代について一定水準の年金額を確保するという財政方式のもとで運営されている。

  このような財政方式のもとでは、当初の見込みに比べて少子高齢化が進行すると、
  高齢者の相対的な増加によって年金給付が相対的に増加することから、
  年金財政の給付と負担の均衡を保つためには、現役世代の負担の増加、
  又は年金受給者の給付の抑制が必要となる。

  このため、平成 16 年改正において（　　B　　）が導入されたが、
  どの程度給付水準を調整する必要があるかについては、
  \begin{itemize}
    \item 高齢化や少子化がどの程度まで進行するか
    \item 女性や高齢者の労働市場への参加がどの程度進み、年金制度の支え手がどの程度増加するか
    \item 経済成長がどの程度達成され、賃金の伸びや積立金の運用収入がどの程度見込まれるか
  \end{itemize}
  等、現在及び将来の人口や経済の動向に左右されるものである。
  
  そこで、年金の財政状況を定期的に確認するために、少なくとも5年ごとに「財政検証」を実施し、
  おおむね（　　C　　）年という長期の財政収支の見通し、
  （　　B　　）の開始及び終了年度の見通し並びに給付水準の見通しを作成して、財政
  状況を検証することとされたのである。 また、財政検証においては、
  将来の人口や経済について一定の前提を設定する
  が、将来は不確実であり、その検証を行う時点において使用可能なデータを用い
  最善の努力を払ったとしても実績と前提に乖離が生じるのは避けられない。

  このため、財政検証では、時間の経過とともに新たに蓄積された最新のデータ
  を用いて諸前提を設定し直した上で、現実の軌道を出発点として新たな出発点からおおむね
  （　　C　　）年間の見通しを作成する。それとともに、幅のある複数の前提を設定し、
  将来の人口や経済の姿に応じて、年金の財政状況を検証し、
  将来の給付水準等がどのようになるかを示しているのである。

  財政検証の結果は、人口や経済を含めた将来の状況を正確に見通す予測 (forecast)というよりも、
  人口や経済等に関して現時点で得られるデータの将来の年金財政への投影(projection)という
  性格のものであることに留意が必要であり、結果を幅広く解釈した上で、
  今後の進むべき方向を読み解くことが重要である。

\end{fillblank}

\begin{answerbox}
  \begin{enumerate}[A: ]
    \item 賦課方式
    \item マクロ経済スライド
    \item 100
  \end{enumerate}
\end{answerbox}

\begin{explanationbox}
  類題：年金2 2024 問題1(4)
\end{explanationbox}

\begin{fillblank}[数理人試験 2024 問題1 設問7 改題]
  「令和6(2024)年財政検証結果レポート」第３章の「第４節　経済前提の設定　1.経済要素と年金財政」
  に関する記述について、以下の空欄を埋めなさい。
  \\

公的年金の財政の将来見通しを作成するにあたっては、人口に関する前提のほか、
経済状態に対応した報酬の上昇や（　　A　　）、（　　B　　）スライド等を将来推計に織り込むため、
（　　B　　）上昇率、（　　A　　）上昇率、（　　C　　）の 3 つ経済要素を経済前提として設定する必要がある。
\\

\noindent \textbf{(1)収入、支出の計算スキームにおける経済前提の役割}

年金給付の将来見通しの作成に当たって、年金額は毎年度、（　　B　　）、（　　A　　）の変動に
応じて改定されることが基本であるため、年金改定率の将来見通しのため（　　B　　）上昇率、
（　　A　　）上昇率が用いられる。また、2 階部分の報酬比例年金は被保険者であった期間の
（　　D　　）に比例し支給される仕組みであるため、（　　D　　）の将来推計のためにも賃金上昇率が用いられる。

保険料収入の将来見通しについてみると、厚生年金は被保険者の標準報酬の一定割合(（　　E　　）\%)
で保険料が賦課されるため、報酬比例年金と同様に（　　D　　）の将来推計にあたって（　　A　　）上昇率が用いられる。
また、国民年金の保険料は（　　A　　）スライドする仕組みであるため、ここでも（　　A　　）上昇率、
（　　B　　）上昇率が用いられる。

さらに、国庫負担の将来見通しは（　　F　　）の一定割合で計算され、
（　　C　　）を用いて積立金の運用収入の将来見通しが計算されている。
このようにして公的年金の収入、支出の将来見通しが作成される。


（以下略）

\end{fillblank}

\begin{answerbox}
  \begin{enumerate}[A: ]
    \item 賃金
    \item 物価
    \item 運用利回り
    \item 標準報酬
    \item 18.3
    \item 年金給付
  \end{enumerate}
\end{answerbox}



\begin{explanationbox}
  5年前と全く同じ文章
\end{explanationbox}

\begin{fillblank}[数理人試験 2025 問題3 設問1]
  次は、「「国民年金及び厚生年金に係る財政の現況及び見通しの関連試算－令和６（2024）年オプ
ション試算結果－」の記載の一部である。以下の空欄を埋めなさい。
\\

オプション試算の内容
\begin{enumerate}[1. ]
  \item \underline{被用者保険の更なる適用拡大}
  \begin{enumerate}[(1)]
    \item 被用者保険の適用対象となる企業規模要件の廃止と5人以上個人事業所に係る非適用業種の
解消を行う場合 （約90万人）\\
    （略）
    \item (1)に加え、短時間労働者の賃金要件の撤廃又は（　　A　　）の引上げにより同等の効果が
得られる場合 （約200万人）\\
    （略）
    \item (2)に加え、5人未満の個人事業所も適用事業所とする場合 （約270万人）\\
    （略）
    \item 所定労働時間が週（　　B　　）の全ての被用者を適用する場合 （約860万人）\\
    （略）
  \end{enumerate}
  \item \underline{基礎年金の拠出期間延長・給付増額}
  
  基礎年金の保険料拠出期間を現行の40年（20～59歳）から45年（20～64歳）に延長し、拠出期間
が伸びた分に合わせて基礎年金が増額する仕組みとした場合
  \item \underline{マクロ経済スライドの調整期間の一致}
  
  基礎年金（１階）と報酬比例部分（２階）に係るマクロ経済スライドの調整期間を一致させた場合
  \item \underline{（　　C　　）}
  
  就労し、一定以上の賃金を得ている65歳以上の老齢厚生年金受給者を対象に、当該老齢厚生年金の
一部または全部の支給を停止する仕組み （（　　C　　））の見直しを行った場合
  \item \underline{標準報酬月額の上昇}
  
  厚生年金の標準報酬月額の上限（現行（　　D　　））の見直しを行った場合
\end{enumerate}
\end{fillblank}

\begin{answerbox}
  \begin{enumerate}[A: ]
    \item 最低賃金
    \item 10時間以上
    \item 在職老齢年金制度
    \item 65万円
  \end{enumerate}
\end{answerbox}

\begin{comment}
  

\begin{fillblank}[自作 2026/11/4]
  以下の表は、給付水準調整終了後のモデル年金の所得代替率の見通し及び調整終了年度を各試算別にまとめたものである。
  以下の空欄を埋めなさい。

  \[
    足下の所得代替率：（　　\text{A}　　）％
  \]
  \begin{table}[htbp]
    \centering
    \begin{tabular}{|c|c|c|c|} \hline 
      \multirow{2}{*}{試算内容} &   \multicolumn{3}{|c|}{経済前提} \\ 
      \cline{2-4} & 高成長実現ケース & 成長型経済移行・継続ケース & 過去30年投影ケース \\ \hline
      現行の仕組み & （　　B　　）(2039) & （　　C　　）(2037) & （　　D　　）(2057) \\ \hline
      \begin{tabular}{l}
        1. 被用者保険の更なる適用拡大\\
        (1)企業要件廃止\&5人以上\\
           \;\;\;\; 個人事業所の非適用業種の廃止 
      \end{tabular} & - & 58.6\%(2035) & 51.3\%(2054)\\ \hline
      \begin{tabular}{l}
        1. 被用者保険の更なる適用拡大\\
        (2)(1)\&賃金要件撤廃
      \end{tabular} & - & 59.3\%(2034) & 51.8\%(2052)\\ \hline
      \begin{tabular}{l}
        1. 被用者保険の更なる適用拡大\\
        (3)(2)\&5人未満の個人事業所も\\
         \;\;\;\; 適用事業所に
      \end{tabular} & - & 60.7\%(2028) & 53.1\%(2048)\\ \hline
      \begin{tabular}{l}
        1. 被用者保険の更なる適用拡大\\
        (4)(3)\&週10時間以上に拡大
      \end{tabular} & - & 61.2\%(調整なし) & 53.1\%(2038)\\ \hline
    2. 基礎年金の拠出期間延長・給付増額 & - & 64.7\%(2038) & 57.3\%(2055)\\ \hline
    3. マクロ経済スライドの調整期間の一致 & - & 61.2\%(調整なし) & 56.2\%(2036)\\ \hline
    \end{tabular}

  \end{table}
  
\end{fillblank}

\begin{answerbox}
  \begin{enumerate}[A: ]
    \item 61.2\%
    \item 56.9\%
    \item 57.6\%
    \item 50.4\%
  \end{enumerate}
\end{answerbox}
\end{comment}


\newpage